%% kapa: exam=1 round=37 year=2026 period=1 form=A subject=real_estate subject_ko=부동산학원론 seq=28 number=108
%% source: Q-Net 감정평가사 기출문제 (cst00309) artlSeq=5258924 file=raw/1008/5258924/2026년 제37회 감정평가사 1차 1교시.pdf
%% conversion: kordoc 4.24.51 -> tools/gen1.py (auto); review=proofread-sonnet; flags=none
\begin{question}{108}
이자율과 할인율이 연 8 \%로 일정할 때, A, B, C를 크기 순서로 나열한 것은? (단, 주어진 자료에 한하며, 모든 현금흐름은 연말에 발생함)
\begin{framedtext}
○ A: 2차년도부터 6차년도까지 매년 250만원씩 받는 연금의 현재가치 \par
○ B: 2차년도부터 6차년도까지 매년 200만원씩 받는 연금의 6차년도의 미래가치 \par
○ C: 1차년도에 40만원을 받고 매년 전년대비 3 \%씩 수령액이 증가하는 성장형 영구연금의 현재가치 \par
○ 연금의 현가계수(8 \%, 5년): 3.9927, 연금의 현가계수(8 \%, 6년): 4.6229 \par
○ 연금의 내가계수(8 \%, 5년): 5.8666, 연금의 내가계수(8 \%, 6년): 7.3359 \par
\end{framedtext}
\begin{choices}
  \item A > B > C
  \item A > C > B
  \item B > A > C
  \item C > A > B
  \item C > B > A
\end{choices}
\end{question}
\answer{3}
